%% ============================================================
%%  Capítulo 7 – Integridad de Señal
%% ============================================================
\chapter{Requisitos de Integridad de Señal}
\label{chap:si}

\section{Metodología de Identificación}
\label{sec:si_metodologia}

Las señales críticas del sistema ALMA fueron identificadas a partir de tres criterios:
(1)~frecuencia de operación y tiempos de subida/bajada de la señal;
(2)~sensibilidad a la perturbación (señales analógicas de audio HD);
(3)~consecuencia del fallo (pérdida de imagen, pérdida de audio, violación de
privacidad). Para cada señal crítica se especifica el riesgo que presenta y la
restricción concreta que deberá respetarse en el layout del Entregable~4.

% ── Alta velocidad ────────────────────────────────────────────
\section{Señales de Alta Velocidad}
\label{sec:si_alta_velocidad}

\subsection{Interfaz MIPI DSI}
\label{subsec:si_mipi}

La interfaz MIPI DSI conecta el SoC A311D2 con el panel táctil a color, operando
a velocidades de varios Gbps. Es la señal de mayor frecuencia del sistema y la
más sensible a discontinuidades de impedancia \reqref{REQ-G-01}.

\begin{sibox}[SI-01: MIPI DSI -- Impedancia Diferencial y Length Matching]
  \begin{itemize}
    \item \textbf{Impedancia diferencial:} 100\,$\Omega$~$\pm$~10\,\%.
      Implementación: par microstrip diferencial en capa~1, plano de referencia
      GND en capa~2 continuo bajo toda la ruta.
    \item \textbf{Length matching entre carriles:} tolerancia máxima de
      $\pm$5\,mm entre los 4~carriles de datos. Verificar con la regla
      \textit{Matched Lengths} del DRC de Altium.
    \item \textbf{Skew intra-par:} máximo 5\,ps por par diferencial.
    \item \textbf{Longitud máxima de ruta:} $\leq$100\,mm para mantener la
      atenuación dentro del presupuesto del receptor del panel.
    \item \textbf{Prohibiciones:} no pasar vías en la ruta del MIPI DSI;
      no cruzar sobre cortes del plano de GND; clearance mínimo de 3$\times$W
      respecto a otras señales (regla de los 3W).
    \item \textbf{Terminación:} AC coupling (condensadores de serie de 100\,nF)
      según especificación MIPI Alliance \cite{mipi_dsi}.
  \end{itemize}
\end{sibox}

\subsection{USB 2.0 (Depuración / Conectividad)}
\label{subsec:si_usb}

\begin{sibox}[SI-02: USB 2.0 -- Par Diferencial]
  \begin{itemize}
    \item \textbf{Impedancia diferencial:} 90\,$\Omega$~$\pm$~10\,\%.
    \item \textbf{Length matching intra-par:} $\leq$0.5\,mm de diferencia
      entre D+ y D--.
    \item \textbf{Restricciones de vías:} minimizar el número de vías en el
      par diferencial; si son inevitables, usar vías pequeñas y añadir una vía
      de GND adyacente para mantener el retorno.
    \item \textbf{Plano de referencia:} continuo bajo toda la longitud del par.
  \end{itemize}
\end{sibox}

% ── Audio ─────────────────────────────────────────────────────
\section{Señales de Audio y Analógicas}
\label{sec:si_audio}

\subsection{Bus I2S -- Audio HD 24-bit/96\,kHz}
\label{subsec:si_i2s}

El reloj de bit (BCLK) del bus I2S para audio de 24-bit a 96\,kHz opera a:
\[
  f_{BCLK} = 2 \times 24 \times 96\,000 = 4.608\,\text{MHz}
\]
Esta frecuencia no impone restricciones severas de SI en términos de integridad
de tiempo de propagación, pero la calidad de audio de 24-bit sí exige un piso de
ruido bajo en las señales del bus \reqref{REQ-F-01}.

\begin{sibox}[SI-03: I2S -- Integridad de Reloj y Aislamiento de Ruido]
  \begin{itemize}
    \item \textbf{Length matching:} BCLK, LRCLK y DATA deben tener longitudes
      similares ($\pm$10\,mm) para mantener el alineamiento temporal de las muestras.
    \item \textbf{Plano de referencia:} GND continuo bajo todas las pistas I2S;
      no cruzar sobre cortes del plano ni sobre la zona de los reguladores Buck.
    \item \textbf{Aislamiento de zona:} las pistas I2S deben mantenerse alejadas
      de los nodos de conmutación (SW) de los reguladores NB679/NB680 mediante
      una separación mínima de 2\,mm o una fila de vías de GND como blindaje.
    \item \textbf{Terminación:} las líneas I2S son punto a punto (SoC -- códec),
      sin necesidad de terminación adicional a esta frecuencia.
  \end{itemize}
\end{sibox}

\subsection{Señales PDM de Micrófonos MEMS}
\label{subsec:si_pdm}

Las líneas PDM llevan el reloj y los datos digitales de los micrófonos MEMS.
El ruido en la línea de reloj PDM degrada directamente la calidad de la captura
de voz \reqref{REQ-F-02}.

\begin{sibox}[SI-04: PDM -- Aislamiento del Ruido del Altavoz]
  \begin{itemize}
    \item \textbf{Separación mecánica:} los micrófonos MEMS deben ubicarse
      en la PCB lo más alejado posible del amplificador Clase~D y del altavoz,
      tanto eléctrica como mecánicamente, para evitar retroalimentación acústica
      y acoplamiento de las corrientes de conmutación del amplificador.
    \item \textbf{Ruteo:} las pistas PDM deben evitar cruzar sobre o cerca de las
      pistas de potencia del amplificador (VCC5, pistas de salida al altavoz).
    \item \textbf{Blindaje:} se recomienda una fila de vías de GND perimetral a
      lo largo de las pistas PDM en la zona de mayor exposición al ruido.
    \item \textbf{Plano de referencia:} GND continuo bajo las pistas PDM.
  \end{itemize}
\end{sibox}

\subsection{Ruta de Audio Analógica -- Códec hacia Amplificador}
\label{subsec:si_audio_analog}

La señal analógica entre la salida del códec y la entrada del amplificador Clase~D
es la más sensible al ruido de todo el sistema. Cualquier interferencia que se
acople en esta ruta será amplificada y reproducida por el altavoz.

\begin{sibox}[SI-05: Audio Analógico -- Zona y Plano de Referencia]
  \begin{itemize}
    \item \textbf{Zona analógica definida:} el códec y la ruta analógica deben
      situarse en un área delimitada de la PCB, separada de los circuitos digitales
      y de potencia.
    \item \textbf{Plano GND analógico:} se recomienda una región del plano GND
      dedicada a la zona analógica, conectada al resto del plano GND en un único
      punto (conexión de estrella) para evitar que las corrientes de retorno digitales
      crucen bajo la zona analógica.
    \item \textbf{Pistas cortas:} la longitud de la ruta analógica (códec -- amp)
      debe minimizarse. Ubicar ambos componentes adyacentes en el layout.
    \item \textbf{Prohibición:} ninguna pista de señal digital o de potencia
      switching debe cruzar perpendicular ni paralela a las pistas de señal analógica
      en capas adyacentes.
  \end{itemize}
\end{sibox}

% ── Control ───────────────────────────────────────────────────
\section{Señales de Control}
\label{sec:si_control}

\subsection{Buses I2C}
\label{subsec:si_i2c}

\begin{sibox}[SI-06: I2C -- Capacitancia de Bus y Resistencias de Pull-up]
  Los buses I2C (uno para el códec, uno para el controlador táctil) son de baja
  velocidad ($\leq$400\,kHz). La restricción principal es la capacitancia de bus:
  \begin{itemize}
    \item \textbf{Capacitancia máxima de bus:} 400\,pF según especificación I2C.
      Con pistas de 0.1\,mm sobre FR-4, la capacitancia es $\approx$0.5\,pF/mm,
      lo que permite longitudes de hasta 800\,mm -- sin restricción práctica en
      esta PCB.
    \item \textbf{Resistencias de pull-up:} 4.7\,k$\Omega$ a 3.3\,V para
      velocidad estándar (100\,kHz) o 2.2\,k$\Omega$ para modo fast (400\,kHz).
      Ubicar los resistores de pull-up cerca del master (SoC) en la PCB.
  \end{itemize}
\end{sibox}

\subsection{Señal de Silenciado Físico (SPSF)}
\label{subsec:si_mute}

\begin{sibox}[SI-07: Hard-wired Mute -- Continuidad Garantizada]
  La señal de silenciado físico (\reqref{REQ-NF-01}) es una interconexión
  hard-wired entre el switch mecánico DPST y la línea de alimentación o reloj
  de los micrófonos PDM. Sus requisitos de SI son:
  \begin{itemize}
    \item \textbf{Continuidad del circuito:} verificar en el ERC de Altium que
      la red de los micrófonos PDM muestra una apertura en el switch (pines del
      componente Switch conectados en serie en la red).
    \item \textbf{Independencia del plano digital:} las pistas del circuito de
      silenciado y del LED indicador no deben compartir retorno con las pistas
      de señal del SoC. Se recomienda un retorno directo a GND independiente.
  \end{itemize}
\end{sibox}

% ── PDN ───────────────────────────────────────────────────────
\section{Integridad de Potencia (PDN)}
\label{sec:si_pdn}

\subsection{Estrategia de Desacoples}

\begin{sibox}[SI-08: Desacoples -- Criterio de Ubicación y Valores]
  \begin{itemize}
    \item \textbf{Filosofía:} cada pin de alimentación de cada CI debe tener un
      condensador de desacoplo cerámico (100\,nF, X7R) ubicado a $\leq$1\,mm
      del pin. Los CIs de mayor consumo (SoC, amplificador) deben tener además
      un condensador bulk ($\geq$10\,$\mu$F).
    \item \textbf{Ubicación en Altium:} los condensadores de desacoplo deben
      asignarse al mismo net que el pin de alimentación al que sirven y colocarse
      en el lado del componente donde están los pines de VCC.
    \item \textbf{Conexión:} la vía del condensador de desacoplo al plano GND
      debe ser la más corta posible. Nunca interponer una pista larga entre el
      condensador y el plano.
  \end{itemize}
\end{sibox}

\subsection{Plano de GND y Continuidad de Retorno}

\begin{sibox}[SI-09: Plano GND -- Sin Cortes bajo Señales Críticas]
  \begin{itemize}
    \item El plano de GND de la capa~2 de la PCB-CTRL debe ser continuo,
      sin ranuras ni cortes bajo ninguna pista de señal de alta velocidad
      (MIPI DSI, I2S, PDM, USB).
    \item \textbf{Regla de retorno:} la corriente de retorno sigue el camino
      de mínima inductancia, que es directamente bajo la pista de señal.
      Un corte en el plano fuerza el retorno a rodear el corte, creando un
      lazo de área mayor y emisiones de radiación aumentadas.
    \item \textbf{Inspección en Altium:} verificar con la función
      \textit{Highlight Net} $>$ GND + \textit{Polygon Pour} que no existen
      zonas desconectadas del plano de GND bajo las rutas críticas.
  \end{itemize}
\end{sibox}

% ── Señales inter-PCB ─────────────────────────────────────────
\section{Señales en la Interconexión entre PCBs}
\label{sec:si_inter_pcb}

Como se indicó en el Capítulo~\ref{chap:particion}, las únicas señales que cruzan
entre la PCB-PWR y la PCB-CTRL son rieles de alimentación DC. No hay señales de
alta velocidad en la interconexión inter-PCB. Esto elimina los problemas de
discontinuidad de impedancia y retorno de corriente de alta frecuencia que
normalmente introduce un conector.

La única consideración de SI en la interconexión es la caída de tensión resistiva
en los pines del conector bajo corriente DC, que debe minimizarse seleccionando
un conector de baja resistencia de contacto y distribuyendo la corriente de cada
riel entre múltiples pines en paralelo.

% ── Tabla resumen ─────────────────────────────────────────────
\section{Tabla Resumen de Restricciones de Integridad de Señal}
\label{sec:si_tabla}

\begin{table}[H]
  \centering
  \caption{Resumen de restricciones de integridad de señal del sistema ALMA.}
  \label{tab:si_resumen}
  \begin{tabularx}{\textwidth}{C{1.2cm} L{2.5cm} C{1.5cm} L{5.5cm} L{3cm}}
    \toprule
    \rowcolor{udea-green}
    \textcolor{white}{\textbf{ID}} &
    \textcolor{white}{\textbf{Señal}} &
    \textcolor{white}{\textbf{PCB}} &
    \textcolor{white}{\textbf{Restricción principal}} &
    \textcolor{white}{\textbf{Evidencia en Altium}} \\
    \midrule
    SI-01 & MIPI DSI  & CTRL & 100\,$\Omega$ diff., length match $\pm$5\,mm, no vías & Matched Lengths DRC \\
    \rowcolor{row-alt}
    SI-02 & USB 2.0   & CTRL & 90\,$\Omega$ diff., match intra-par $\leq$0.5\,mm & Diff. Pair Rule DRC \\
    SI-03 & I2S       & CTRL & GND continuo, separación de Buck, match $\pm$10\,mm & Visual inspection \\
    \rowcolor{row-alt}
    SI-04 & PDM Mics  & CTRL & Alejado del amp., fila de vías GND perimetral & Visual inspection \\
    SI-05 & Audio analógico & CTRL & Zona analógica, GND estrella, pistas cortas & Visual inspection \\
    \rowcolor{row-alt}
    SI-06 & I2C       & CTRL & Pull-up correcto, cap. bus $\leq$400\,pF & ERC check \\
    SI-07 & Mute HW   & CTRL & Continuidad, retorno GND independiente & ERC check \\
    \rowcolor{row-alt}
    SI-08 & PDN -- Desacoples & CTRL & 100\,nF a $\leq$1\,mm de cada pin VCC & Visual inspection \\
    SI-09 & Plano GND & CTRL & Sin cortes bajo señales críticas & Polygon DRC \\
    \rowcolor{row-alt}
    SI-PWR-01 & PDN PWR  & PWR & $\Delta$V $\leq$3\,\% en rieles críticos & Cálculo (Anexo B) \\
    \bottomrule
  \end{tabularx}
\end{table}
